\documentclass[12pt]{article}

\usepackage{amsmath}
\usepackage{geometry}
\geometry{margin=1in}

\setlength{\parskip}{0.5em}

\begin{document}

\title{Grade 11 Physics -- Circular Dynamics Practice 1 \\ \large{Horizontal Circular Motion and Centripetal Force}}
\author{}
\date{}
\maketitle

\vspace{1em}

\textbf{Instructions:}

\begin{enumerate}
\item Draw a Free-Body Diagram (FBD) for every problem before doing calculations.
\item Show all steps clearly, starting from Newton's Second Law ($\vec{F}_{\text{net}} = m\vec{a}_c$).
\item Include units in every step and in the final answer. Use $g = 9.8 \, \text{m/s}^2$ unless specified otherwise.
\end{enumerate}

\noindent\rule{\textwidth}{0.4pt}

\vspace{1em}

\section*{Question 1: Identifying the Centripetal Force}

For each of the following scenarios of uniform circular motion, state which physical force (e.g., Tension, Gravity, Normal Force, Static Friction) acts as the centripetal force $F_c$.

\begin{itemize}
\item[(a)] A car making a left turn on a flat, horizontal road: \underline{\hspace{4cm}}
\vspace{0.8cm}

\item[(b)] The Moon orbiting the Earth: \underline{\hspace{4cm}}
\vspace{0.8cm}

\item[(c)] A child standing against the inner wall of a spinning "Rotor" ride (the floor drops out, but the child doesn't slide down): \underline{\hspace{4cm}}
\vspace{0.8cm}

\item[(d)] A puck tied to a string being swung in a horizontal circle on a frictionless ice rink: \underline{\hspace{4cm}}
\end{itemize}

\vspace{1cm}

\section*{Question 2: Period, Frequency, and Acceleration}

A $0.200 \, \text{kg}$ rubber stopper is swung in a horizontal circle on the end of a $1.20 \, \text{m}$ string. It completes $15$ full revolutions in $10.0 \, \text{s}$. (Assume the string remains perfectly horizontal for this problem).

\begin{itemize}
\item[(a)] Calculate the frequency ($f$) and the period ($T$) of the rubber stopper.
\vspace{5cm}

\item[(b)] Calculate the speed ($v$) of the rubber stopper.
\vspace{5cm}

\item[(c)] Calculate the centripetal acceleration ($a_c$) of the stopper.
\vspace{6cm}

\end{itemize}

\newpage

\section*{Question 3: Mass on a String (Frictionless Table)}

A block of mass $m_1 = 1.50 \, \text{kg}$ is on a frictionless horizontal table. It is attached to a string that passes through a small hole in the center of the table. The block is moving in a uniform circle of radius $r = 0.800 \, \text{m}$ at a constant speed of $4.00 \, \text{m/s}$.

\begin{itemize}
\item[(a)] Draw a Top-View and Side-View Free-Body Diagram for the sliding block.
\vspace{5cm}

\item[(b)] What force provides the centripetal force in this situation?
\vspace{6cm}

\item[(c)] Calculate the tension in the string required to keep the block moving in this circle.
\vspace{6cm}

\end{itemize}

\newpage

\section*{Question 4: Car on a Flat Curve}

A $1200 \, \text{kg}$ car enters a flat, unbanked curve with a radius of $65.0 \, \text{m}$. The coefficient of static friction ($\mu_s$) between the car's tires and the dry asphalt is $0.800$. 

\begin{itemize}
\item[(a)] Draw the Free-Body Diagram for the car (rear-view perspective). 
\vspace{6cm}

\item[(b)] Calculate the maximum safe speed at which the car can negotiate the curve without skidding.
\vspace{8cm}

\end{itemize}

\newpage

\section*{Question 5: The "Rotor" Ride (Challenge)}

In an amusement park "Rotor" ride, people stand against the inner wall of a large hollow cylinder that is spinning. Once the cylinder spins fast enough, the floor drops away, but the people do not slide down the wall. 

The radius of the cylinder is $4.50 \, \text{m}$. The coefficient of static friction ($\mu_s$) between a person's clothes and the wall is $0.400$.

\begin{itemize}
\item[(a)] Draw the Free-Body Diagram for a person against the wall. Hint: The wall pushes horizontally \textit{inward} on the person. Which force prevents the person from falling down?
\vspace{6cm}

\item[(b)] Write Newton's Second Law equations for both the vertical and horizontal directions.
\vspace{6cm}

\item[(c)] Determine the minimum speed $v$ the person must be moving at so they do not slide down the wall.
\vspace{6cm}

\end{itemize}

\end{document}
